% report_use.tex — Complete usage template for hustreport.cls
% Compile with XeLaTeX (or LuaLaTeX)

\documentclass[linux]{hustreport} % options: win | mac | linux | fonts | thesis

% -----------------------------------------------------------------
% Basic metadata (used by cover and PDF metadata)
% -----------------------------------------------------------------
\title{社会调查报告题目}
\author{张三}
\date{\today}

\dept{管理学院}
\major{信息管理与信息系统 2022 级}
\advisor{李四}
\team{王五、赵六}
\studentID{2022000000}
\reportID{报告编号：HUST-2026-001}

% -----------------------------------------------------------------
% Task book content (optional)
% -----------------------------------------------------------------
\renewcommand{\taskTargetText}{调查目标内容……}
\renewcommand{\taskContentText}{方案设计内容……（可很长，会自动跨页）}
\renewcommand{\taskExpectedText}{预期成果内容……}
\renewcommand{\taskScheduleText}{计划进度内容……}

\begin{document}

% ----------------------------
% Cover + Task Book
% ----------------------------
\makecover
\maketaskbook

% ----------------------------
% Abstracts (optional)
% ----------------------------
\begin{cnabstract}{关键词1；关键词2；关键词3}
这里写中文摘要……
\end{cnabstract}

\begin{enabstract}{keyword1; keyword2; keyword3}
Write English abstract here…
\end{enabstract}

% ----------------------------
% Table of Contents
% ----------------------------
\tableofcontents
\newpage

% ----------------------------
% Main content
% ----------------------------
\section{引言}
本节介绍研究背景与意义。

\section{方法}
\subsection{数据采集}
描述数据来源与采集方法。

\subsection{分析方法}
说明分析框架与指标体系。

\section{结果}
\subsection{关键发现}
呈现主要调查结果。

\section{结论}
总结与建议。

% ----------------------------
% Figures and tables examples
% ----------------------------
\begin{figurerow}{不同指标对比}[2]
    \figitem{figure/py2.png}{指标 A}
    \figitem{figure/py3.png}{指标 B}
\end{figurerow}

\begin{table}[H]
    \centering
    \caption{示例表格}
    \begin{tabular}{ccc}
        \toprule
        项目 & 数值 & 备注 \\
        \midrule
        A & 1.23 & 示例 \\
        B & 4.56 & 示例 \\
        \bottomrule
    \end{tabular}
\end{table}

% ----------------------------
% Code block examples
% ----------------------------
\begin{codeblock}[python]{示例代码}
print("hello")
\end{codeblock}

% \filecode[python]{外部代码示例}{path/to/code.py}

% ----------------------------
% Notes / review examples (optional)
% ----------------------------
\todo{待补充：完善调研样本的描述}
\fixme{此处数据需核对}
\note{这是备注}

\review{需要标注的文本}{侧边批注内容}
\review[style=highlight,color=blue,side=below]{高亮文本}{行间批注}

% ----------------------------
% Appendix (optional)
% ----------------------------
\begin{appendixblock}
\section{附录标题}
附录内容……
\end{appendixblock}

% ----------------------------
% Bibliography
% ----------------------------
% Method A: use .bib file
% \MyBibliography{refs}

% Method B: manual bibliography
\begin{thebibliography}{99}
\bibitem{ref1} 作者. 标题. 期刊, 年份.
\end{thebibliography}

% ----------------------------
% Acknowledgement (optional)
% ----------------------------
\begin{acknowledgement}
感谢……
\end{acknowledgement}

\end{document}
